\chapter*{Abstract}
\begin{spacing}{1.5}
	
	Depression is a critical worldwide mental health problem that reduces quality of life and increases the risk of suicide. Traditional methods like interviews and questionnaires are limited and lack scalability and objectivity. Social media's vast user-generated content offers a unique opportunity for advancing depression detection methodologies. So this study presents a deep learning framework to identify and classify depression severity (non-depressed, mild, moderate, severe) from social media text. 
	
	This approach leverages advanced natural language processing techniques to analyze tweets and other user-generated content. The method begins with rigorous preprocessing—normalizing, tokenizing, and cleaning raw text—followed by feature extraction using a pre-trained BERT model. These contextual embeddings are then refined by a hierarchical attention network (HAMN) that captures both fine-grained (word-level) and broader (sentence-level) patterns in the data. The final model classifies texts into four distinct categories ranging from non-depressed to severely depressed. 
	
	Experimental results on a DEPTWEET dataset of over 40,000 tweets show promising performance. The model achieves 83\% accuracy, with precision at 80\%, recall at 82\%, and F1-score at 81\%, excelling at non-depressed cases but struggling with severity levels. Results show transformer-based deep learning is effective for early depression detection and timely mental health intervention. Future work will focus on further improving model interpretability and integrating additional data modalities to enhance diagnostic accuracy.
	
\end{spacing}